\section{Inversão de dados em eletrorresistividade}
\label{subsec:inversao_de_dados_em_er}

A Teoria da inversão de dados é utilizada em eletrorresistividade para estimar modelos geológicos que justifiquem os valores das medidas físicas adquiridas durante levantamentos efetuados na superfície. Os dados coletados no campo podem ser expressas através de um vetor de dados $d = {[~d_1,~d_2,~d_3~...~d_n~]}^{T}$, onde $T$ significa transposta. Da mesma forma, os parâmetros do modelo podem ser representados como os elementos de um vetor modelo $m = [~\rho_1,~\rho_2,~\rho_3~...~\rho_n,~E_1,~E_2,~E_3~...~E_n-1]^{T}$. O parâmetro do modelo ($2n - 1$) é a última camada, sendo o semi-espaço. Os parâmetros $\rho_n$ e $E_n$ são os valores de resistividade e espessura, respectivamente.

Esses vetores são relacionados através de um operador de modelagem $g$, conforme equação abaixo \cite{koefoed1979resistivity}:
\begin{align}
    \label{eq:dgm}
    d=g(m)
\end{align}

No processo de inversão, estimamos um modelo $m_{est}$ que minimiza a função objetivo ou função \textit{fitness}, a qual é uma medida de discrepância entre os dados observados e dados calculados. Conforme a equação abaixo:
\begin{equation}
    \label{eq:fit}
    F=\frac{1}{N}|| d^{obs}-d^{cal}||^2
\end{equation}


Os dados observados ($d_{obs}$) e os dados calculados ($d_{calc}$) bem como o número total de dados ($N$) são calculados, para dados sintéticos e reais, usando a equação \ref{eq:fit}. \cite{Benjumea_1d_2021}

Na inversão deve ser utilizado uma curva de campo para decidir o número de camadas e resistividades aparente e definir um limite de pesquisa de valores de resistividade e espessura da camada. \citeonline{Sen1993} utilizam uma abordagem que serve como direcionamento para o processo de inversão desse trabalho.